\section{Conclusioni}\label{sec:con}

Il lavoro risponde a una domanda: che cosa può dichiarare un sistema di percezione nell'intervallo in cui un ostacolo nasconde il soggetto. La risposta adottata è un'uscita che tiene distinti l'osservato, il previsto con la sua regione e il suo livello, e il non noto, e che non dichiara mai «libero» ciò che non è stato osservato. I risultati hanno tre livelli di prova diversi, che qui si tengono separati.

\paragraph{Che cosa è dimostrato, nel modello dichiarato \prova{D}.}
\begin{itemize}
\item Nessuna cella in ombra, fuori campo, oltre portata o letta per meno di metà dell'impronta è dichiarata «libero» (per costruzione); e, se la profondità valida è corretta, la posa è esatta e il suolo è piano, i raggi dei pixel che leggono il suolo di una cella «libero» sono privi di superfici fino al suolo, a meno del margine di lettura, all'istante d'acquisizione (Proposizione~\ref{prop:map-sicurezza}). L'enunciato non copre la cella intera né la colonna sopra di essa: la sicurezza della cella intera e la tenuta con un errore d'inclinazione entro $0{,}5^\circ$ sono verificate in simulazione su 300 condizioni dichiarate, non dimostrate \prova{S} (\S\ref{sec:map-banco}). Fuori da queste ipotesi l'enunciato cade: i controesempi sono il vetro, l'ostacolo che non rimanda segnale e l'errore d'inclinazione di $1^\circ$ (\S\ref{sec:map-controesempi}).
\item Un soggetto nascosto, con velocità non superiore a $\vmax$, resta nell'insieme di possibile presenza $\Pres(t)$, che si riduce solo per osservazione e mai per il tempo trascorso (Proposizioni~\ref{prop:occ-contenimento} e~\ref{prop:occ-monotonia}): a posa esatta, o con il libero eroso del margine di posa, su suolo piano e con la lettura per punto; con la mappa a celle il contenimento poggia sulla verifica in simulazione \prova{S}.
\item La fascia di guardia $w$ contiene il soggetto se vale la condizione~\eqref{eq:occ-fascia}, l'ostacolo è fermo e il soggetto resta nel campo e in portata; con $w=0{,}40\unita{m}$ la velocità coperta, per calcolo, è $1{,}13\unita{m/s}$, e $0{,}56\unita{m/s}$ se il bordo dell'ombra sbaglia di $0{,}20\unita{m}$ (Proposizione~\ref{prop:occ-fascia}) \prova{C}.
\item La regione di previsione è contenuta nell'area d'avviso e, con l'ostacolo fermo e l'ombra stimata costante, non si svuota finché la traccia è nell'ombra; per ogni ingresso riconosciuto, l'area d'avviso con la previsione contiene quella senza (Proposizione~\ref{prop:occ-annidamento}). Un'etichetta di classe non produce un «libero» e, per ogni ingresso riconosciuto, non riduce l'avviso (Proposizione~\ref{prop:occ-etichetta}).
\item La massa non è identificabile da misure di posizione, e l'accelerazione non è stimabile con precisione utile alla frequenza e al rumore assunti (Proposizione~\ref{prop:sti-massa} per la massa \prova{D}; conto del \S\ref{sec:sti-massa} per l'accelerazione \prova{C}).
\item La geometria d'osservazione è in forma chiusa: visibilità del suolo, ombra di un ostacolo, durata dell'occlusione, rumore di profondità ed effetto di un errore di posa (\S\ref{sec:geo}).
\end{itemize}

\paragraph{Che cosa è verificato in simulazione \prova{S}.}
\begin{itemize}
\item \emph{Con modello ideale del sensore.} La regione è emessa a ogni ciclo dell'occlusione: 0 passaggi con un'interruzione su $23\,402$, da 6000 soggetti distinti; la premessa, cioè l'ingresso nell'ombra riconosciuto, manca in 928 passaggi su $23\,402$ ($3{,}97\,\%$), e in 24 ($0{,}10\,\%$) con l'avviso acceso anche per le tracce chiuse entro $0{,}75\unita{m}$ da un'ombra, regola attivata solo dopo la prova d'ingresso sul suolo vuoto (\S\ref{sec:lim-salute}). Con un ostacolo largo $2\unita{m}$, il bordo dell'ombra noto e il soggetto sul proprio percorso, la regione contiene il soggetto con estremo inferiore dell'intervallo sopra $0{,}90$ fino a $2\unita{s}$ e sopra $0{,}85$ a $3\unita{s}$, in 30 combinazioni su 30. Con un pannello largo $1\unita{m}$ la soglia di $0{,}90$ non è raggiunta a $1\unita{s}$ ($942/1045$, $\ell^-=0{,}882$) e oltre non si dichiara alcuna cifra (\S\ref{sec:val}).
\item La regola del «libero» della cella non dichiara libera alcuna cella sopra un ostacolo in 300 condizioni dichiarate, con errore d'inclinazione entro $0{,}5^\circ$; lo zero vale per quell'insieme di condizioni e non è una frequenza (\S\ref{sec:map}).
\item \emph{Su profondità simulata dalle schede tecniche.} Un soggetto nascosto entro $3\unita{m}$ non resta senza avviso (0 fotogrammi su 4324 a 424 pixel di larghezza e 0 su 4325 a 848, e 0 su 3608 per risoluzione in una replica). L'identità alla ricomparsa, che con la configurazione di riferimento tiene in 0 passaggi su 90, con i gruppi a celle forti e le condizioni aggiuntive del riaggancio tiene in 176 passaggi su 180 a 424 pixel di larghezza e in 174 su 180 a 848, in un dominio ristretto: un adulto solo, a $0{,}5\unita{m}$ dietro il pannello, entro $2{,}6\unita{m}$, rumore nominale (\S\ref{sec:sti}, \S\ref{sec:val}). Una cella «libero» dentro l'ombra vera di un ostacolo: 0 su $2\,907\,032$ celle\,$\cdot$\,fotogramma in una prima serie di 135 sequenze, 5 su $3\,992\,186$ in una seconda di 381 (una al livello di rumore nominale); entrambe con la variante che conserva il «libero» per tre fotogrammi, non con la regola della cella adottata. Non è risolto se le 5 dipendano dalla regola o dal bordo della maschera di riferimento.
\item \emph{Su traiettorie pubbliche di pedoni} \prova{P}. Il modello di moto a velocità quasi costante, con $q=0{,}2\unita{m^2/s^3}$, copre il soggetto in $0{,}977$, $1{,}000$ e $1{,}000$ dei casi a 1, 2 e $3\unita{s}$ (219 tratti di 9 persone; numerosità efficace fra 9 e 219; posizioni note al millimetro fino all'ultimo istante, senza profondità): riguarda il modello di moto, non la catena.
\item Il costo di calcolo è in media $6{,}35\unita{ms}$ per ciclo su profondità simulata, su un filo di un processore da portatile, e $81\unita{ms}$ su una registrazione pubblica affollata (\S\ref{sec:cos}).
\end{itemize}

\paragraph{Dove le simulazioni mostrano che il sistema cede.} Chi svolta o esce dal lato lontano dell'ombra è nella regione, a 2 e $3\unita{s}$, in una quota di casi fra $0{,}57$ e $0{,}86$ nelle prove con il corpo esteso e fra $0{,}13$ e $0{,}70$ in quelle con il soggetto puntiforme; chi parte da fermo accanto all'ostacolo, in $0{,}11$--$0{,}26$ dei casi fra 1 e $3\unita{s}$ nelle prove di base; con il bordo dell'ombra stimato e non riconosciuto la quota a $1\unita{s}$ scende a $0{,}30$--$0{,}46$ (estremo basso: chi prosegue). In questi casi non si dichiara un livello e resta la sola area d'avviso. Un secondo soggetto mai visto che esce dalla stessa ombra riceve l'identità del primo: 2400 casi su 2400 nelle prove di base con modello ideale del sensore, fra $0{,}91$ e $0{,}995$ nelle prove d'insieme con modello ideale, 19 su 20 (76--99\,\%) su profondità simulata dalle schede tecniche; nessuna delle regole provate lo ripara, e per questo l'identità è presunta e l'avviso non ne dipende. Su profondità simulata dalle schede tecniche un pannello fermo è dichiarato «in movimento» in 79 passaggi su 90 con i gruppi per componenti connesse; con i gruppi a celle forti in 0 su 140 nel dominio dell'identità al ritorno e in 7 su 10 fuori da esso. Dietro un ostacolo fermo l'avviso non si spegne.

\paragraph{Che cosa resta non verificato sperimentalmente \prova{N}.}
\begin{itemize}
\item \emph{Ogni grandezza su un sensore fisico.} Nessuna prova è stata eseguita su una camera: rumore, pixel senza dato, errore del bordo dell'ombra, copertura della regione, identità alla ricomparsa, latenza e costo sulla scheda di bordo non sono misurati. La latenza non è misurata: la somma delle voci note dà $155\unita{ms}$ in media, o circa $131\unita{ms}$ con un altro dato del costruttore per il tratto del sensore; le due fonti non concordano e nessuna voce è misurata sulla catena.
\item \emph{La continuità su profondità simulata.} La regola della regione persistente è stata eseguita solo con modello ideale del sensore (\S\ref{sec:lim-limiti}).
\item \emph{Il termine dell'avviso per osservazione.} Che l'avviso non termini senza osservazione vale per costruzione; che termini quando il suolo è rivisto dopo la rimozione dell'ostacolo non è sostenuto: nella realizzazione provata resta acceso dopo $2\unita{s}$ in 10 prove su 10 (72--100\,\%), per una causa non determinata.
\item \emph{Il classificatore.} È un componente richiesto dal sistema; in questo lavoro non è realizzato né valutato: è simulato con etichette assegnate a mano. Nessuna classe diversa da «persona» è stata simulata nella catena completa, e la regola che lega la classe alla fascia di guardia è stata provata sulla sola fascia (0 riduzioni dell'avviso su $20\,000$ sequenze di etichette, due volte). La sicurezza delle indicazioni non dipende dall'etichetta.
\item \emph{L'avviso tempestivo per classe e per portatore.} Con $3\unita{m}$ di portata, sensore fermo e un tempo di reazione ipotizzato di $1{,}0\unita{s}$, l'avviso è tempestivo per calcolo fino a $1{,}44\unita{m/s}$: 3 classi su 15 alla velocità tipica, nessuna alla massima. Il sensore in moto è trattato solo come analisi dei requisiti.
\item \emph{L'uso da parte di persone.} Lo schermo e i segnali sonori non sono stati provati con persone, vedenti o no.
\end{itemize}

Le misure che deciderebbero ciascuno di questi punti, con la soglia che li smentirebbe, sono elencate in \S\ref{sec:lim-misure}. Il sistema svolge solo osservazione, monitoraggio e avviso: non è un dispositivo di protezione.
